\paragraph{\trHomotopy。}\label{sub:homotopy}
考察 $\Ob\scr{Q}$（其中 $\scr{Q}=\mathit{nu}\text{-}\mathit{fun}(\scr{A},\scr{B})$）中的两个%
\trNonunital{}\trAFunctor{} $\scr{F}_0$ 与 $\scr{F}_1$，它们在\trObject{}上的作用相同。令
$D=\scr{F}_0-\scr{F}_1\in\hom^1_{\scr{Q}}(\scr{F}_0,\scr{F}_1)$ 为由 $D^0=0$、
$D^d=\scr{F}_0^d-\scr{F}_1^d$（$d>0$）所定义的预\trNaturalTransformation。它其实是一个真正的%
\trNaturalTransformation，即 $\mu^1_{\scr{Q}}(D)=0$；若 $D=\mu^1_{\scr{Q}}(T)$，其中
$T\in\hom^0_{\scr{Q}}(\scr{F}_0,\scr{F}_1)$ 还满足 $T^0=0$，就称 $\scr{F}_0$ 与
$\scr{F}_1$ 是\trHomotopic{}的。至少从表面上看，这个名称是有依据的：两个\trHomotopic{}的%
\trFunctor{}在\trCohomology{}上诱导同一个\trFunctor，即 $H(\scr{F}_0)=H(\scr{F}_1)$。
\trHomotopy{}在与一个固定的\trFunctor{}作左、右\trComposition{}之下保持。这就是说，若 $T$ 是从
$\scr{F}_0$ 到 $\scr{F}_1$ 的一个\trHomotopy，则 $\scr{R}^1_{\scr{G}}(T)$ 是
$\scr{F}_0\circ\scr{G}$ 与 $\scr{F}_1\circ\scr{G}$ 之间的一个\trHomotopy；而
$\scr{L}^1_{\scr{G}}(T)$ 则是 $\scr{G}\circ\scr{F}_0$ 与 $\scr{G}\circ\scr{F}_1$ 之间的一个%
\trHomotopy。

\trHomotopy{}是一个\trEquivalence{}关系。原则上这并不难验证，但计算可能显得有点令人困惑，原因
如下。由于我们所考虑的\trFunctor{} $\scr{F}_k$ 在\trObject{}上的作用全都相同，\trGraded%
\trVectorSpace{} $\hom_{\scr{Q}}(\scr{F}_k,\scr{F}_l)$ 对一切 $k$、$l$ 都是同一个，
而我们将会想把预\trNaturalTransformation{}从一个搬到另一个那里去。为免混淆，我们将采用记号
$T_{[kl]}$，用以表明 $T$ 被看作落在 $\hom_{\scr{Q}}(\scr{F}_k,\scr{F}_l)$ 之中。在这一
术语下，最初步的观察是：对任意三个\trFunctor{} $\scr{F}_j$、$\scr{F}_k$、$\scr{F}_l$，有
\begin{equation}\label{eq:homotopy-basic-observation}
\begin{split}
\mu^1_{\scr{Q}}(T_{[jl]})
&=\mu^1_{\scr{Q}}(T_{[jk]})-\mu^2_{\scr{Q}}(\scr{F}_k-\scr{F}_l,T_{[jk]})\\
&=\mu^1_{\scr{Q}}(T_{[kl]})+\mu^2_{\scr{Q}}(T_{[kl]},\scr{F}_j-\scr{F}_k).
\end{split}
\end{equation}

为证传递性，设 $T_1$ 是从 $\scr{F}_0$ 到 $\scr{F}_1$ 的一个\trHomotopy，而 $T_2$ 是从
$\scr{F}_1$ 到 $\scr{F}_2$ 的一个\trHomotopy。利用\eqref{eq:homotopy-basic-observation}
可以看出，
\begin{equation}\label{eq:transitivity-homotopy}
T=(T_1)_{[02]}+(T_2)_{[02]}+\mu^2_{\scr{Q}}(T_2,T_1)
\end{equation}
满足
\begin{equation*}
\begin{split}
\mu^1_{\scr{Q}}(T)&=\scr{F}_1-\scr{F}_2-\mu^2_{\scr{Q}}(\scr{F}_1-\scr{F}_2,T_1)\\
&\hspace{1.5em}+\scr{F}_0-\scr{F}_1+\mu^2_{\scr{Q}}(T_2,\scr{F}_0-\scr{F}_1)
+\mu^1_{\scr{Q}}\mu^2_{\scr{Q}}(T_2,T_1)=\scr{F}_0-\scr{F}_2,
\end{split}
\end{equation*}
因而是从 $\scr{F}_0$ 到 $\scr{F}_2$ 的一个\trHomotopy。对称性可循同样的思路：设 $T_1$ 是从
$\scr{F}_0$ 到 $\scr{F}_1$ 的一个\trHomotopy。考虑映射
\begin{equation}\label{eq:symmetry-maps-phi-psi}
\begin{aligned}
\phi\colon\hom_{\scr{Q}}(\scr{F}_1,\scr{F}_0)&\longrightarrow
\hom_{\scr{Q}}(\scr{F}_1,\scr{F}_1),&
\phi(T)&=(-1)^{|T|}\mu^2_{\scr{Q}}(T_1,T)+T_{[11]},\\
\psi\colon\hom_{\scr{Q}}(\scr{F}_1,\scr{F}_1)&\longrightarrow
\hom_{\scr{Q}}(\scr{F}_0,\scr{F}_1),&
\psi(T)&=\mu^2_{\scr{Q}}(T,T_1)+T_{[01]}.
\end{aligned}
\end{equation}
由于每个公式里都有最后一项，它们都是\trIsomorphism；而且应用
\eqref{eq:homotopy-basic-observation} 可以看出，它们是\trChainMap。最后，再用一次
\eqref{eq:homotopy-basic-observation} 就得到
\begin{equation*}
\begin{aligned}
\phi\bigl((\scr{F}_0-\scr{F}_1)_{[10]}\bigr)&=\mu^1_{\scr{Q}}\bigl((T_1)_{[11]}\bigr),\\
\psi\bigl(\mu^1_{\scr{Q}}((T_1)_{[11]})\bigr)&=\mu^1_{\scr{Q}}\mu^2_{\scr{Q}}
\bigl((T_1)_{[11]},T_1\bigr)+(\scr{F}_0-\scr{F}_1).
\end{aligned}
\end{equation*}
结论是：$-T_1-\mu^2_{\scr{Q}}\bigl((T_1)_{[11]},T_1\bigr)$ 在 $\psi\circ\phi$ 下的\trPreimage，
就是从 $\scr{F}_1$ 到 $\scr{F}_0$ 的一个\trChainHomotopy。

\begin{remark}\label{thm:homotopy-rational-interpretation}
上面这些计算乃至一般\trHomotopy{}的含义，当 $\scr{B}$ 具有严格的\trIdentityMorphism{}时会变得
清楚一些；见下文 (2c)。与此同时，这里给出另一种解释，它来自有理\trHomotopy{}理论
[68, Chapter X]。考虑如下的 dg 代数：$I=\mathbb{K}u_0\oplus\mathbb{K}u_1\oplus\mathbb{K}h$，
其中 $|u_0|=|u_1|=0$，$|h|=1$；其\trDifferential{}与乘法（略去因\trDegree{}原因而为零的
表达式）为 $\partial u_0=h=-\partial u_1$，以及
\begin{equation*}
\begin{aligned}
u_0u_0&=u_0,& u_1u_1&=u_1,& u_0u_1&=u_1u_0=0,\\
u_1h&=h,& hu_1&=u_0h=0,& hu_0&=h.
\end{aligned}
\end{equation*}
$H(I)\cong\mathbb{K}$，以 $u_0+u_1$ 为\trGenerator，于是投射到由 $u_k$ 生成的子空间上的那些映射
$f_k\colon I\to\mathbb{K}$（$k=0$，$1$）都是\trQuasiisomorphism。若 $\scr{B}$ 是一个%
\trNonunital{}\trAinfty-代数，那么\trTensorProduct{} $I\otimes\scr{B}$ 可以直截了当地定义出来：
在底层\trGraded\trVectorSpace{}的\trTensorProduct{}上，令
$\mu^1(i\otimes b)=i\otimes\mu^1_{\scr{B}}(b)+(-1)^{|i|+|b|}\partial(i)\otimes b$；而对
$d\ge 2$，令 $\mu^d(i_d\otimes b_d,\ldots,i_1\otimes b_1)=(-1)^{\lhd}\,i_d\cdots i_1\otimes
\mu^d_{\scr{B}}(b_d,\ldots,b_1)$，其中 $\lhd=\sum_{j>k}(|b_j|-1)|i_k|$。一个简单的计算
表明：$T$ 是两个 \trAinfty-\trHomomorphism{} $\scr{F}_0,\scr{F}_1\colon\scr{A}\to\scr{B}$ 之间的%
\trHomotopy，当且仅当
\begin{equation*}
\begin{split}
\scr{G}(a_d,\ldots,a_1)
&=u_0\otimes\scr{F}_0(a_d,\ldots,a_1)+u_1\otimes\scr{F}_1(a_d,\ldots,a_1)\\
&\hspace{2em}+(-1)^{\clubsuit_d}\,h\otimes T(a_d,\ldots,a_1)
\end{split}
\end{equation*}
定义出一个 \trAinfty-\trHomomorphism{} $\scr{G}\colon\scr{A}\to I\otimes\scr{B}$。整个情形可以
总结成如下的交换图
\begin{tikzcd}
& & & \scr{B} \\
\scr{A} \ar[rr,"\scr{G}"]
  \ar[urrr,bend left=30,"\scr{F}_0"]
  \ar[drrr,bend right=30,"\scr{F}_1"']
& & I\otimes\scr{B} \ar[ru,"f_0\otimes\id"] \ar[rd,"f_1\otimes\id"] & \\
& & & \scr{B}.
\end{tikzcd}
同样的做法对 \trACategory{}与\trFunctor{}也适用，只需作显然的修改。
\end{remark}
